% Lösungen Einheit 4 von 5 – Dezimalzahlen (zusammenbau v0.1)
\einheitenkopf{Einheit 4 von 5 · Dezimalzahlen}

%% TODO Lösung der Erklärzeile 21g) fehlt in der Bank
\erg{21}{a) $2$ \quad b) $0{,}6$ \quad c) $0{,}47$ \quad d) $\frac{8}{10}$ \quad e) $\frac{39}{100}$ \quad f) $0{,}83$ \quad g) \ldots \quad h) $0{,}07$ \quad i) $0{,}4$ am vierten, $1{,}3$ am dreizehnten Teilstrich nach der $0$ \quad j) $\frac{2}{5} = \frac{4}{10} = 0{,}4$ \quad k) $\frac{6}{10} = \frac{3}{5}$ \quad l) $7 : 8 = 0{,}875$ \quad m) $4 : 9 = 0{,}444\ldots = 0{,}\overline{4}$ \quad n) $\frac{13}{5} > \frac{5}{2}$, denn $\frac{13}{5} = 2{,}6$ und $\frac{5}{2} = 2{,}5$; $\sqrt{5} \approx 2{,}24$}

21i)

\zahlenstrahl[xmin=0,xmax=2,xstep=0.1,karo=0.6]{0/0, 1/1, 2/2, 0.4/0{,}4, 1.3/1{,}3}

\erg{22}{a) $2{,}047$}
\erg{23}{a) $4{,}3$ und $4{,}4$}
\erg{24}{a) Sie hat den Bruchstrich als Komma gelesen. Richtig: $\frac{1}{4} = \frac{25}{100} = 0{,}25$. \quad b) Erweitern mit $5$ ergibt $\frac{45}{100}$, und $45$ Hundertstel schreibt man $0{,}45$. \quad c) $0{,}2$ \quad d) $\frac{3}{5}$ l $= \frac{6}{10}$ l $= 0{,}6$ l}
